\section{Scope and Limitations}
The current model is intended for architecture-stage exploration. It does not perform detailed routing, extraction, static timing analysis, IR/EM analysis, thermal simulation, clock-tree synthesis, or signoff DRC. Process scaling from 28 nm toward advanced nodes is a calibrated abstraction assembled from public information and user-editable assumptions; it is not a substitute for a foundry PDK. Likewise, PHY depth, SRAM macro overhead, defect density, and wafer cost can dominate an actual product and must be replaced with project-specific data when available.

The Nangate45 experiment validates only the reported topology families, library, placement conditions, and congestion projection. Its correlations must not be transferred automatically to TSMC, Samsung, Intel, SMIC, GF, Huawei Tao, W2W stacks, or new CPU templates. CPU reference entries separate public architectural anchors from estimated physical budgets. Where cache or power data are uncertain, the UI exposes the assumption rather than presenting it as vendor measurement.

Future work should calibrate the connection model against multiple open PDKs, add obstacle-aware coarse routing, evaluate partition objectives using real netlists, and establish per-template synthesis baselines. A second direction is hierarchical multi-die exploration in which fabric, bonding-site, package escape, and thermal constraints are optimized jointly rather than independently.
